% Run GCT_KRON_20261120; GPT-6 Astra Pro; effort not exposed.
\section{The positive algebra also survives on every proper partial determinant}

\subsection*{A-036. A polynomial lower bound on every proper Hessian ray}
Let \(n\geq3\), \(R=n^2\), and \(3\leq r<R\).
Use the partial-Hessian notation of A-035:
\[
 q_r=r(n-2),\qquad D_r=r+1,\qquad
 \lambda_{n,r}=(n+q_r+1,3,2^{r-2}).
\]
Let \(F_i\) be the normalized binary coefficients of \(F\),
and \(H_i\) those of the \(r\)-variable Hessian
\(H_r(F)\). Define the centered binary polynomials
\(I_2,\ldots,I_n\) exactly as in A-034, and set
\[
 Z_r=F_0H_1-F_1H_0,\qquad U_i=H_0^i I_i
                                     \quad(2\leq i\leq n).
                                                               \tag{A-036.1}
\]
Then these \(n\) functions are algebraically independent
in the determinant coordinate ring:
\[
 \boxed{\quad
 \mathbb C[Z_r,U_2,\ldots,U_n]
       \lhook\joinrel\longrightarrow
 \bigoplus_{s\geq0}
      \mathbb C[\mathcal D_n]_{sD_r}^{\,U}[s\lambda_{n,r}].
 \quad}                                                        \tag{A-036.2}
\]
Their ray degrees are \(1,2,\ldots,n\). In particular,
\[
 \boxed{\quad
 K_n(s\lambda_{n,r})
       \geq p_n(s):=[z^s]\prod_{i=1}^n(1-z^i)^{-1}
       \qquad(s\geq0,\ 3\leq r<n^2).
 \quad}                                                        \tag{A-036.3}
\]
This is a lower bound, not an exact formula for the
entire proper-partial determinant multiplicity.
At \(s=1\), A-035 gives the exact value one.

The bound contrasts with the full-variable case \(r=R\),
where A-034 gives the smaller exact determinant value
\(h_n(s)=[z^s]\prod_{i=2}^n(1-z^i)^{-1}\).
In the full case the bracket \(Z_R\) vanishes on the
determinant; in every proper dimension at least three,
it supplies an algebraically independent degree-one
generator alongside the same centered construction.
The restriction \(r\geq3\) is essential to this proof
and statement.

\paragraph{An explicit isotropic direction at an invertible matrix.}
Let \(P=\det_n\), and equip \(W=\operatorname{Mat}_n(\mathbb C)\)
with its nondegenerate Hessian bilinear form at the identity,
\[
             Q(C,D)=\operatorname{tr}(C)\operatorname{tr}(D)
                                          -\operatorname{tr}(CD).
                                                               \tag{A-036.4}
\]
Take
\[
 v=E_{12}+E_{21}+E_{13}-E_{31},\qquad
 B_0=\operatorname{diag}(0,1,-1,0,\ldots,0).
                                                               \tag{A-036.5}
\]
Then
\[
 \operatorname{tr}v=0,\qquad \operatorname{tr}(v^2)=0,\qquad
 Q(v,I)=Q(v,B)=0
 \]
for every diagonal matrix \(B=\operatorname{diag}(b_1,\ldots,b_n)\).
Thus \(v\) is a nonzero isotropic vector, and every
diagonal matrix lies in \(v^\perp\).
The plane \(\langle I,B_0\rangle\) is nondegenerate:
its \(Q\)-matrix is
\[
                         \begin{pmatrix}n(n-1)&0\\0&-2\end{pmatrix}.
                                                               \tag{A-036.6}
\]
The form on \(v^\perp/\mathbb Cv\) is nondegenerate
of dimension \(R-2\). Within this quotient, take the
orthogonal complement of \(\langle I,B_0\rangle\);
it is nondegenerate of dimension \(R-4\).
Choose a nondegenerate subspace of dimension \(r-3\)
there and lift a basis to fixed vectors
\(C_4,\ldots,C_r\in v^\perp\), orthogonal to
\(\langle I,B_0\rangle\).
This is possible because
\[
                              r-3\leq R-4.
\]
When \(r=3\), there are no extra vectors.

For variable diagonal \(B\), consider the ordered frame
\[
                         (I,B,v,C_4,\ldots,C_r).
                                                               \tag{A-036.7}
\]
Let \(\Delta(B)\) be the determinant of the \(Q\)-Gram
matrix on the ordered vectors \(I,B,C_4,\ldots,C_r\).
At \(B=B_0\), (A-036.6) and the choice of extras show
\(\Delta(B_0)\ne0\).
Hence, for \(B\) in a nonempty open subset of the
diagonal matrices, (A-036.7) is an independent
\(r\)-frame and the restricted Hessian has radical
exactly \(\mathbb Cv\).
The vector \(v\) is orthogonal to every vector in
the frame for every diagonal \(B\), so its Gram
determinant vanishes identically on this entire
irreducible family, not just at \(B_0\).

\paragraph{The Hessian vanishes simply along a direction visible to the bracket.}
For a general ordered \(r\)-frame
\(\mathbf A=(A_1,\ldots,A_r)\), write
\[
 F_{\mathbf A}(x_1,\ldots,x_r)
                   =P\left(\sum_{j=1}^r x_j A_j\right).
 \]
View \(F_i,H_i,Z_r,I_i\) as polynomial functions of
the entries of \(\mathbf A\).
On the family (A-036.7), \(F_0=P(I)=1\) and \(H_0=0\).
We compute \(H_1\) there.

Differentiating the determinant Hessian at \(I\) in
the direction \(B\), and using the trace identities
above, gives
\[
       \left.\frac{d}{dt}\right|_{t=0}
              \operatorname{Hess}P_{I+tB}(v,v)
             =2\operatorname{tr}(Bv^2)=2(b_2-b_3).
                                                               \tag{A-036.8}
\]
For example, differentiating
\[
 P(X)\{\operatorname{tr}(X^{-1}v)^2
                         -\operatorname{tr}(X^{-1}vX^{-1}v)\}
\]
at \(X=I\) proves this directly: the terms containing
\(\operatorname{tr}v\) or \(\operatorname{tr}(v^2)\)
are zero.
Since the row and column corresponding to \(v\)
are zero at \(t=0\), only their diagonal cofactor
contributes to the derivative of the restricted
Gram determinant. Thus
\[
                  q_r H_1=2(b_2-b_3)\Delta(B).
                                                               \tag{A-036.9}
\]
At \(B_0\) this is nonzero, and therefore
\[
                  Z_r=H_1\ne0
          \quad\hbox{on a nonempty open subset of (A-036.7)}.
                                                               \tag{A-036.10}
\]

The parameter derivative in this argument has the
precise normalization needed below.
Changing the first frame vector from \(A_1\) to
\(A_1+tA_2\), while holding the others fixed, is
the variable shear \(x_2\mapsto x_2+t x_1\) on
\(F_{\mathbf A}\). Its determinant is one.
Partial-Hessian covariance therefore gives
\[
 \left.\frac{d}{dt}\right|_{t=0}
        H_0(A_1+tA_2,A_2,\ldots,A_r)=q_r H_1(\mathbf A).
                                                               \tag{A-036.11}
\]
Thus the derivative used to detect simple vanishing
is a genuine derivative of the polynomial \(H_0\)
on the full affine frame-parameter space.

\paragraph{The centered binary invariants remain independent on the zero family.}
On (A-036.7), the binary restriction is
\[
                 F(x,y)=\det(xI+yB)=\prod_{j=1}^n(x+b_jy).
 \]
Its raw coefficients \(c_1,\ldots,c_n\) are the
elementary symmetric polynomials in \(b_1,\ldots,b_n\);
their map is dominant over \(\mathbb C\).
The normalized coefficients are
\(F_i=c_i/\binom ni\), with \(F_0=1\).
As in A-034, the map
\[
       (F_1,\ldots,F_n)\longmapsto(F_1,I_2,\ldots,I_n)
 \]
is triangular with all diagonal entries one.
Therefore \(I_2,\ldots,I_n\) are algebraically
independent on the family (A-036.7).
Removing the proper closed subset where
\(\Delta(B)(b_2-b_3)=0\) does not change this fact.

\paragraph{A simple prime factor in the full frame ring.}
Let
\[
                     \mathcal A=\mathbb C[(A_j)_{ab}]
 \]
be the polynomial ring of all \(r\)-frame parameters;
no independence condition is imposed in its definition.
The polynomial \(H_0\in\mathcal A\) is nonzero,
because (A-036.9)--(A-036.11) give a nonzero derivative.
Its specialization to (A-036.7) is zero identically.
Factor \(H_0\) into irreducibles in \(\mathcal A\).
Since the coordinate ring \(\mathbb C[b_1,\ldots,b_n]\)
of the family is a domain, some irreducible factor
\(h\) specializes identically to zero there.

That factor has multiplicity one, and no second
irreducible factor can vanish identically on the
same family. Otherwise every first derivative
of \(H_0\) would specialize to zero, contradicting
(A-036.9)--(A-036.11).
Moreover (A-036.10) gives \(h\nmid Z_r\).
For every nonzero polynomial \(T\), the preceding
binary independence gives
\[
                     h\nmid T(I_2,\ldots,I_n).
 \]
Consequently the prime valuation in
\(\operatorname{Frac}(\mathcal A)\) satisfies
\[
 \operatorname{ord}_h H_0=1,\qquad
 \operatorname{ord}_h Z_r=0,\qquad
 \operatorname{ord}_h T(I_2,\ldots,I_n)=0\quad(T\ne0).
                                                               \tag{A-036.12}
\]
The family itself need not be a divisor.
It locates a simple prime factor in the full frame ring,
which is all that the valuation proof requires.

\paragraph{Independence in the actual determinant ring.}
The partial covariance calculation of A-035 gives
\[
 \deg Z_r=D_r,\quad \operatorname{wt}Z_r=\lambda_{n,r},
 \qquad
 \deg U_i=iD_r,\quad \operatorname{wt}U_i=i\lambda_{n,r}.
 \]
Suppose a nonzero polynomial relation among these
functions vanished on the determinant orbit closure.
Use the coefficient-degree grading to choose a
weighted-homogeneous relation of ray degree \(s\),
and write it as
\[
                       \sum_j Z_r^j T_j(U_2,\ldots,U_n)=0,
 \]
where \(T_j\) has weighted degree \(s-j\).
Every independent \(r\)-frame extends to a basis of
\(\operatorname{Mat}_n\), so its form \(F_{\mathbf A}\)
is the restriction of a determinant orbit translate.
Thus the proposed relation pulls back to zero on
the dense open set of independent frames, and hence
is an identity in \(\mathcal A\).
After substituting \(U_i=H_0^i I_i\), it reads
\[
                     \sum_j Z_r^j H_0^{s-j}T_j(I_2,\ldots,I_n)=0.
 \]
For every nonzero block \(T_j\), (A-036.12) gives
valuation \(s-j\). These values are distinct.
The block with largest \(j\) cannot cancel with
any other block, a contradiction.
This proves the injection (A-036.2), and its
positive monomial basis proves (A-036.3).

\paragraph{What this excludes under internal padding.}
For \(3\leq r<n^2\), the span of the constructed
monomials of ray degree \(s\) has dimension \(p_n(s)\)
already on the determinant. Restriction of the same
universal span to any permanent orbit closure has
dimension at most \(p_n(s)\). Thus this entire
polynomial positive-algebra construction cannot
certify a strict permanent multiplicity advantage
at a proper partial ray, at any stretch.
This conclusion concerns its own highest-weight
vectors, not all highest-weight vectors of that
type.

For any internal padding \(m<n\), every useful
partial dimension satisfies \(r\leq m^2<n^2\),
so the preceding exclusion applies whenever
\(r\geq3\). If \(r>m^2\), the padded permanent
multiplicity vanishes by length. If \(r=2\),
every stretched weight is two-row, and O-006 gives
the full ambient determinant multiplicity, so
again no actual obstruction occurs there.
The additional whole-coefficient-algebra injection
in A-035, for \(r\leq2n\), is stronger than the
particular polynomial-subalgebra injection here.
Neither result excludes further copies outside
the specified algebra at general proper partial
types.

\paragraph{Geometric dictionary and limits.}
The frame, Gram determinant, shear and centered
coefficient constructions are actual polynomial
maps, with rational normalized coefficients.
They have the comodule and torus-weight dictionary
of SGA 3, I, Propositions 4.7.2--4.7.3,
pp.~34--36. The characteristic-zero Schur
interpretation uses SGA 6, VI, Theorem 3.3,
pp.~392--393, and Macdonald, I, Appendix
A(5.3)--(5.6), pp.~155--156, as in A-034--A-035.
The nondegenerate trace pairing, isotropic-frame
construction and prime-factor differentiation are
used over \(\mathbb C\). No claim about arbitrary
special fibers is made.
The proof produces a polynomial subalgebra of the
actual determinant highest-weight ring; it neither
computes that entire ring in the proper dimensions
nor supplies a general padded multiplicity comparison.
